% Generated draft: requires mathematical and source review.
\documentclass[11pt]{article}
\usepackage[margin=25mm]{geometry}
\usepackage{fontspec}
\setmainfont{Latin Modern Roman}
\usepackage{amsmath,amssymb,mathtools,unicode-math}
\setmathfont{Latin Modern Math}
\usepackage{enumitem,needspace,xurl,hyperref}
\setlength{\parindent}{0pt}
\setlength{\parskip}{5pt}
\setlength{\emergencystretch}{3em}
\newcommand{\N}{\mathbb{N}}
\newcommand{\Z}{\mathbb{Z}}
\newcommand{\Q}{\mathbb{Q}}
\newcommand{\R}{\mathbb{R}}
\newcommand{\C}{\mathbb{C}}
\newcommand{\F}{\mathbb{F}}
\newcommand{\A}{\mathbb{A}}
\newcommand{\PP}{\mathbb{P}}
\newcommand{\E}{\mathbb{E}}
\newcommand{\eps}{\varepsilon}
\newcommand{\dd}{\,\mathrm{d}}
\begin{document}
\section*{UnsolvedMath ID: 30004427}
\textbf{Source problem code:} catalogue entry
\subsection{Definitions and mathematical statement}
The supplied record concerns a pair \((M,D)\) described as a maximally degenerate positive Calabi--Yau pair admitting a weak section. It does not specify the definitions of these hypotheses, the ambient geometric category, or the meaning of a weak section. These hypotheses must therefore be understood in the intended conventions of the cited report; they cannot be made self-contained from the supplied record alone.

Write \(\mathcal W(X)\) for the wrapped Fukaya category of the intended symplectic space \(X\), and write \(A:=SH^0(X)\) for its degree-zero symplectic cohomology algebra. The record does not define \(X\), its relationship to \((M,D)\), its symplectic structure, or the coefficient and grading conventions. It names functors \(\pi^*\) and \(\pi_*\), but supplies neither their constructions nor their source and target categories.

\textbf{Conjecture (supplied formulation).} For every maximally degenerate positive Calabi--Yau pair \((M,D)\) admitting a weak section, the associated functors \(\pi^*\) and \(\pi_*\) make \(\mathcal W(X)\) a categorical crepant resolution of \(A=SH^0(X)\). In particular, the conjecture requires that
\[
\mathcal W(X)\text{ is relatively Calabi--Yau over }SH^0(X).
\]
The meanings of categorical crepant resolution and relative Calabi--Yau structure, including the relevant dimension, the action of \(A\), and the compatibility required of the functors, are not specified in the supplied record. Consequently, this is a schematic statement of the source conjecture, pending recovery of those definitions. The relative Calabi--Yau requirement is presented as a necessary part of the conjecture; no converse is asserted here.
\subsection{Short English statement}
The supplied conjecture asks whether every pair satisfying the report's maximal degeneration, positivity, Calabi--Yau, and weak-section hypotheses yields a wrapped Fukaya category that resolves its degree-zero symplectic cohomology algebra through the indicated functors. A relative Calabi--Yau property is required. Essential geometric and categorical definitions are missing, so this statement-only draft requires human mathematical review.
\subsection{Sources}
\begin{itemize}
\item Supplied local UnsolvedMath registry record, permanent ID \(30004427\), title \emph{Wrapped Fukaya Categories as Crepant Resolutions}, exact source problem code \texttt{catalogue entry}, in the Oberwolfach Reports Open Problems collection. No individual proposer is supplied; the record gives proposed year \(2020\).
\item Supplied report citation: \emph{Mini-Workshop: Superpotentials in Algebra and Geometry} (2021), Oberwolfach Reports, pp.\ 569--600, DOI \texttt{10.4171/owr/2020/10}. Source URL and supplied literature-source link, labeled ``original OWR report'': \url{https://doi.org/10.4171/owr/2020/10}. The record attributes the problem to workshop open-problem discussions documented in this report.
\item Supplied literature assessment, checked on \texttt{2026-08-22}: the record labels the problem open and says that the report reduces the categorical crepant-resolution conjecture to proving a relative Calabi--Yau property, with no general proof of that required property identified in the supplied assessment. This is supplied dated curation, not an independently verified literature assessment. The record explicitly requests independent verification of time-sensitive status and cited claims.
\item The registry labels the statement \texttt{corrected\_verified} and describes editorial repairs as preserving mathematical content. This is supplied registry metadata, not independent verification of the formulation or its mathematical correctness.
\end{itemize}
\end{document}
